\section*{序列空间上的测度}
\phantomsection
\addcontentsline{toc}{section}{序列空间上的测度}

古典概率论中发展最充分的分支之一是极限定理（大数定律、向高斯–拉普拉斯律的趋近，\ldots）；这是对由牵涉极大总体的现象所显示的统计规律性概念的一种深化。这些问题的正确数学表述要求在序列空间上引入测度；这些空间构成有限维空间最明显的推广，是 1920 年前后由弗雷歇、莱维、卢津\ldots{} 开展的"一般分析"研究最钟爱的课题。无论如何，概率论新方法的创立者辛钦与柯尔莫哥洛夫同为卢津的学生，而莱维又很早便转向概率问题，这绝非偶然：这些问题正是新方法的试金石。

序列空间上测度的第一次隐含出现，见于 E. 博雷尔 1909 年论述可数概率的著作{[}32 b{]}。博雷尔一个非常独创的想法，是把他刚刚得到的概率结果应用于证明 0 与 1 之间几乎每个实数的十进展开所具有的性质。这一应用基于下面的基本评注：用每个介于 0 与 1 之间的实数在给定基 \(q\)（\(q \geq 2\)）下展开式中的数字串来定义它；如果随机地依次选取数 \(x\) 的各个数字，彼此独立，并以 \(1/q\) 的等概率取值于 \(0,1,\ldots,q-1\)，则 \(x\) 落入 [0,1] 中某个区间的概率等于该区间的长度。

1923 年，施坦豪斯{[}293{]}严格地建立了这些结果，并描述了博雷尔所考虑的那类无界选择序列的精确数学模型：为简单起见取 \(q = 2\)，并用 \(I\) 记二元集 \(\{0,1\}\)；给 \(I\) 装备测度 \(\mu\)，它由 \(\mu(0) = \mu(1) = \frac{1}{2}\) 定义；乘积空间 \(I^{\mathbf{N}}\) 的元素是取值为 0 或 1 的数构成的序列 \(\varepsilon = (\varepsilon(n))_{n \in \mathbf{N}}\)，而映射 \(\phi : \varepsilon \mapsto \sum_{n \geq 0} \varepsilon(n). 2^{-n-1}\) 除一个可数集外是 \(I^{\mathbf{N}}\) 到区间 \([0,1]\) 上的双射；更有甚者，\(\phi^{-1}\) 把 \([0,1]\) 上的勒贝格测度变为测度 \(P\)，即 \(I^{\mathbf{N}}\) 上各个因子上测度 \(\mu\) 的乘积。事实上，施坦豪斯当时并没有乘积测度的构造；

他利用拟双射 \(\phi\) 的存在性，从 [0,1] 上的勒贝格测度出发构造 \(I^{N}\) 上的测度 P，然后给出 P 的一个公理化刻画。这样得到的同构使人们得以把概率的语言翻译为测度的语言，并应用关于勒贝格积分的已知定理。

在同一篇论文中，施坦豪斯考虑随机级数 \(\sum_{n\geq 0}\sigma_n.a_n\)，其中符号 \(\sigma_{n} = \pm 1\) 是随机选取的，彼此独立且都以同样的概率 \(\frac{1}{2}\) 取值；1928 至 1935 年间，他研究了大量其他的随机级数。另一方面，佩利、维纳与济格蒙德考虑随机傅里叶级数\(^{1}\)，其形如 \(\sum_{-\infty}^{\infty}a_n\exp (2\pi i(nt + \varPhi_n))\)；其中"振幅" \(a_{n}\) 是固定的，而"相位" \(\varPhi_{n}\) 是在 [0, 1] 上均匀分布的独立随机变量。虽然解析上的困难随问题不同而差别极大，但以测度论语言所作的翻译在所有情形都相同，并且代表了博雷尔与施坦豪斯所处理情形的一种推广：问题在于在 \(\mathbf{R}^{\mathbf{N}}\) 上构造一个测度，它是一族测度的乘积，这族测度全都等同于某个质量为 1 的特定正测度 \(\mu\)（在 \(\mathbf{R}\) 上）；例如，上述随机傅里叶级数对应于 \(\mu\) 是 [0, 1] 上勒贝格测度的情形。

为构造这样的乘积测度，可以使用两种方法。第一种是直接方法，由丹尼尔于 1918 年首次建立{[}76 b{]}；1934 年延森重新发现了它{[}173{]}，并对 \(\mu\) 是 \([0,1]\) 上勒贝格测度的情形作了详细研究。第二种方法是寻找与施坦豪斯的技巧类似的技巧，以便化归为 \([0,1]\) 上勒贝格测度的情形；在还没有一般测度论的完整论述可供支配时，这种做法有方便的好处，因为它允许不经新的证明就使用勒贝格的定理。\(^{2}\)

